\documentclass[11pt]{article}

% ---------------------------------------------------------------
% Research direction review: Human-aligned time-series similarity
% Prepared as a learning / discussion document for the supervisor.
% ---------------------------------------------------------------

\usepackage[margin=1in]{geometry}
\usepackage{amsmath,amssymb}
\usepackage{enumitem}
\usepackage{booktabs}
\usepackage{float}
\usepackage{xcolor}
\usepackage[colorlinks=true,linkcolor=blue,citecolor=teal,urlcolor=blue]{hyperref}
\usepackage{tabularx}

\newcommand{\todo}[1]{\textcolor{red}{\textbf{[TODO: #1]}}}
\newcommand{\verify}[1]{\textcolor{orange}{\textbf{[verify: #1]}}}

\title{\textbf{A Human-Aligned, Interpretable Similarity Metric for Time Series}}
\author{Yuli Tshuva}
\date{June 2026}

\begin{document}
\maketitle

\begin{abstract}
\noindent
This document records the proposed framing of the project and a critical assessment
of its novelty and validation strategy. The central thesis is that the proposed
distance --- built from change-point segmentation, hand-crafted shape features, and a
\emph{merge-aware} segment-level dynamic time warping --- is designed not to optimise a
benchmark loss, but to \emph{reproduce the human notion of similarity} between sequences.
The empirical case rests on a human study over \emph{adversarial} examples on which the
proposed metric and standard baselines disagree. We summarise why this framing is strong,
the main threats to the human study's validity, the most closely related prior work
(verified against the literature), and concrete recommendations before final results are
produced.
\end{abstract}

% ===============================================================
\section{The thesis}
% ===============================================================
The goal is a \textbf{human-aligned similarity metric}: a distance that represents the
human notion of similarity between two sequences, rather than one that optimises an
arbitrary objective. The argument has two pillars:
\begin{enumerate}[leftmargin=*]
  \item \textbf{Construction.} The algorithm encodes the structural cues a human reasons
  about --- it segments a sequence at change points, describes each segment with
  interpretable shape features (trend direction/strength, curvature, amplitude, length),
  and aligns segments with a merge-aware DTW that can collapse several segments on either
  side into one matched unit.
  \item \textbf{Validation.} An empirical human study over \emph{adversarial} pairs ---
  examples where the proposed metric and the baselines (DTW, LCSS, Euclidean) rank
  similarity in opposite directions --- in which human raters adjudicate. If humans
  consistently side with the proposed metric, the claim of human-alignment is supported
  by evidence external to the model.
\end{enumerate}

\noindent
Interpretability matters because humans make decisions for reasons they can articulate;
a metric whose alignment is human-readable (a labelled segment-to-segment mapping) is the
mechanism through which the ``reasoning'' can be inspected.

\paragraph{Motivating application.} Comparing signals such as COVID-19 epidemic curves:
locate where a similar trajectory occurred before, and use that precedent to inform the
next decision. This is a query-by-example retrieval setting where \emph{segment-level
shape}, not point-level alignment, is what makes two curves ``the same kind of thing''.

% ===============================================================
\section{Why this framing is strong}
% ===============================================================
\begin{itemize}[leftmargin=*]
  \item It converts the project's main apparent weakness into a strength. The many
  hand-crafted, intuition-driven choices (shape features, trend labels, segment merging)
  would normally be attacked as arbitrary in an ``accuracy'' paper. Under the
  human-alignment thesis they are \emph{defensible by design}: they encode the cues humans
  use.
  \item It moves the work out of the saturated ``yet another DTW variant'' competition
  (where the bar is beating dozens of methods on the UCR archive) into the
  less-occupied niche of \emph{perceptually / human-aligned} similarity.
  \item Interpretability stops being a feature and becomes evidence for the thesis.
\end{itemize}

% ===============================================================
\section{The human study is the paper}
% ===============================================================
With this framing the empirical case rests almost entirely on the human experiment, so
its methodology \emph{is} the contribution. The predictable threats, and how to defuse
them:

\begin{description}[leftmargin=*]
  \item[\textbf{1. Cherry-picking / selection bias (the \#1 threat).}]
  If adversarial examples are hand-selected where the method wins, the rebuttal writes
  itself: ``you chose examples engineered to flatter your method.'' \emph{Mitigation:}
  make selection \emph{mechanical and outcome-blind}. Define ``the baselines disagree with
  us'' by an automatic criterion, sample from that pool, and \textbf{freeze the set before
  observing which side the human took.} Present both directions (ours-high/theirs-low and
  ours-low/theirs-high); this symmetry is what buys credibility.

  \item[\textbf{2. Ask humans the right question.}]
  Do \emph{not} ask people for similarity scores or full rankings --- these are noisy and
  unreliable. Use \textbf{triplets / odd-one-out}: ``Is query $Q$ more similar to $A$ or to
  $B$?'' Each method predicts an answer; agreement with the human majority is the outcome.
  This is the standard robust design for perceptual-similarity studies and sidesteps most
  reliability complaints.

  \item[\textbf{3. Measure inter-rater agreement --- and worry if it is low.}]
  The premise is that ``the human notion of similarity'' is a real, shared target. If
  raters do not agree with \emph{each other} (report Krippendorff's $\alpha$ or similar),
  the target is ill-defined and the thesis weakens. \textbf{Check this early --- it is an
  existential risk, not a detail.}

  \item[\textbf{4. Avoid circularity.}]
  The metric is built from human intuition, then validated against humans. This is sound
  \emph{only} if the validation set is genuinely held out and the metric is never tuned on
  the human judgments. Keep all weight tuning on a separate development set, sealed from the
  study data.

  \item[\textbf{5. Keep a competence floor.}]
  If the only result is ``we win on adversarial cases,'' reviewers suspect a trade.
  Include a sanity check showing the metric is \emph{competitive} (not necessarily SOTA) on
  ordinary cases / a standard benchmark, so the adversarial wins are not bought with
  ordinary-case losses.
\end{description}

\paragraph{Two claims, kept separate.} Conflating these will hurt the paper:
\begin{itemize}[leftmargin=*]
  \item \textbf{Alignment:} the scores match human judgment --- proven \emph{only} by the
  human study.
  \item \textbf{Interpretability:} the method yields a human-readable segment alignment ---
  proven by construction and examples.
\end{itemize}
Interpretability does not imply alignment. Treat them as two distinct contributions, each
with its own evidence.

\paragraph{The application as motivation, not proof.} The COVID retrieval use case is a
strong story for the introduction, but proving it requires showing downstream
\emph{decision} quality, which is a different and harder paper. Safer: a qualitative
query-by-example retrieval demo (``look how sensible the nearest neighbours are'') to
illustrate the application, without resting validity on it.

% ===============================================================
\section{Related work (verified)}
% ===============================================================
The following items were checked against the literature in June 2026. The single most
important one for positioning is DTW+S.

\paragraph{Human perception of time-series similarity.}
Gogolou et al.~\cite{gogolou2019} ran controlled human studies on how well automatic
measures (Euclidean, DTW) match human \emph{perception} of similarity, and showed the
visual representation (line chart vs.\ horizon graph vs.\ colorfield) changes that
perception. This is the must-cite anchor. \textbf{Key differentiation:} they
\emph{measure} how existing metrics align with perception (descriptive); the present work
\emph{constructs} a new metric designed to be human-aligned and interpretable
(constructive). It also primes reviewers to believe ``DTW is already roughly perceptual'',
so the study must beat DTW with statistical significance, not just on hand-picked cases.

\paragraph{Shape-based DTW variants (the methodological neighbours).}
\begin{itemize}[leftmargin=*]
  \item \textbf{shapeDTW} \cite{zhao2018}: encodes a local \emph{shape descriptor} around
  every point and aligns descriptor sequences with DTW, yielding perceptually
  interpretable alignments; beats DTW on 64/84 UCR datasets. Differs from the present work
  in being \emph{point-wise} (one descriptor per timestamp) rather than
  \emph{segment-level}, and has \emph{no human study}.
  \item \textbf{ESDTW} (extrema-based shape DTW)~\cite{esdtw2023}: a more recent shape-DTW
  variant keyed on extrema --- a useful comparison point for the segmentation idea.
\end{itemize}

\paragraph{The closest competitor --- and the open gap.}
\textbf{DTW+S} (Srivastava, 2023)~\cite{dtwpluss} is the nearest neighbour to this project:
it builds an interpretable, ``closeness-preserving'' matrix in which each column is a
\emph{local trend}, then applies DTW --- explicitly to match similar local trends at
similar times, and explicitly motivated by \textbf{clustering epidemic / COVID curves}
(surge $\to$ increase $\to$ peak $\to$ decrease). This overlaps the present project on
three axes at once: shape/local-trend representation, interpretability, and the epidemic
application. \textbf{Crucially, DTW+S contains no human study} --- it is evaluated
computationally (clustering / classification). This is the wedge: the present work's
human-alignment validation is exactly what the closest competitor does not do, and the
segment-level \emph{merge-aware} alignment (many-to-many segment collapsing with a
correlation-gated merge cost) is a mechanism DTW+S does not have. \textbf{This paper must
be cited, contrasted carefully, and ideally used as a baseline.}

\paragraph{Who is already building on DTW+S.}
Checking the papers that cite DTW+S confirms the line of work is active and
epidemic-focused --- and, tellingly, that citers \emph{consume} it as a similarity
primitive rather than scrutinising whether its alignments match human perception (which
is precisely the present project's wedge). The verified citers of DTW+S (the AAAI~2025
published version, ``Aligning time-series by local trends'') are:
\begin{table}[H]
\centering
\renewcommand{\arraystretch}{1.3}
\begin{tabularx}{\textwidth}{p{0.32\textwidth} c c X}
\toprule
\textbf{Citing paper} & \textbf{Year} & \textbf{$\approx$cit.} & \textbf{How it uses DTW+S}\\
\midrule
TrendGNN: Towards Understanding of Epidemics, Beliefs, and Behaviors~\cite{trendgnn}
& 2025 & $<5$ &
Uses DTW+S to build the signal-similarity graph (retaining the top-5 trend-similar edges
per signal) that feeds a GraphSAGE epidemic-forecasting GNN; reports DTW+S-constructed
graphs beat lagged-correlation and random / fully-connected baselines at 2--4 week
horizons. A pure \emph{downstream use} of DTW+S, with no perceptual / human evaluation of
the similarity itself.\\
\bottomrule
\end{tabularx}
\caption{Verified citers of DTW+S.}
\end{table}

\paragraph{Recent baselines and the broader ``human-aligned similarity'' trend.}
\begin{itemize}[leftmargin=*]
  \item \textbf{Multiscale Dubuc}~\cite{dubuc2024}: a 2024 new time-series similarity
  measure --- a fresh baseline to include.
  \item \textbf{Seeing Eye to AI?} (CHI 2025)~\cite{seeingeye2025}: deep-feature metrics
  match human perception of \emph{visualizations} (scatterplots), not time series. Shows
  the ``does this metric see like a human?'' question is active at top HCI venues ---
  useful to cite as evidence the broader direction is timely, while noting it does not
  address time-series shape.
  \item Classic baseline surveys for similarity-measure comparison on
  classification~\cite{empirical2014} remain the reference for the competence-floor
  experiment.
\end{itemize}

\paragraph{Study-design methodology.}
Ordinal / triplet similarity indices~\cite{ordinal2026} formalise interpretable alignment
from \emph{relative} (triplet / quadruplet) comparisons rather than scalar scores --- a
useful reference for justifying the triplet study design and for analysing the collected
judgments.

% ===============================================================
\section{Positioning at a glance}
% ===============================================================
\begin{table}[H]
\centering
\renewcommand{\arraystretch}{1.3}
\begin{tabularx}{\textwidth}{l c c c c X}
\toprule
\textbf{Method} & \textbf{Shape} & \textbf{Segment} & \textbf{Interpret.} & \textbf{Human} & \textbf{Application focus}\\
& \textbf{-based} & \textbf{-level} & & \textbf{study} & \\
\midrule
DTW / LCSS              & no  & no  & low  & no  & general\\
shapeDTW \cite{zhao2018}        & yes & no (point) & med & no  & UCR classification\\
DTW+S \cite{dtwpluss}           & yes & trend cols & yes & \textbf{no} & \textbf{epidemic curves}\\
\textbf{This work}     & yes & \textbf{yes (merge-aware)} & yes & \textbf{yes} & epidemic / decision support\\
\bottomrule
\end{tabularx}
\caption{The combination of segment-level merge-aware alignment \emph{and} a human-alignment
study is the unoccupied cell.}
\end{table}

% ===============================================================
\section{Residual risks reviewers will probe}
% ===============================================================
\begin{itemize}[leftmargin=*]
  \item \textbf{Hyperparameter sensitivity.} Many hand-set constants (feature set,
  hard-coded normalisation means/stds, \texttt{CHANGE\_THRESHOLD}, \texttt{SEGMENT\_PERCENTAGE},
  domain-tuned weights). Ablations / sensitivity analysis are needed to show robustness.
  \item \textbf{Computational cost.} The merge-aware DTW searches over merge spans
  ($d_i, d_j$) and recomputes merged features inside the loop, giving roughly
  $O(n^2 m^2)$ behaviour --- a scalability question on long series. \verify{profile and
  report; consider caching merged features / bounding the max merge span.}
  \item \textbf{``DTW is already perceptual.''} Primed by \cite{gogolou2019}; answered only
  by a statistically significant win in the human study.
\end{itemize}

% ===============================================================
\section{Recommended next steps}
% ===============================================================
\begin{enumerate}[leftmargin=*]
  \item Lock the human-study design \emph{before} producing final results: triplet /
  odd-one-out, outcome-blind mechanical selection of adversarial triplets, pre-registered
  set, inter-rater agreement reported.
  \item Obtain and run \textbf{DTW+S} as a baseline; write the explicit contrast (segment
  merge mechanism + human validation).
  \item Add a competence-floor benchmark (standard datasets) so adversarial wins are not
  read as trades.
  \item Run hyperparameter ablations and a complexity/runtime analysis.
  \item Keep the COVID retrieval as a qualitative motivating demo, not a validity claim.
\end{enumerate}

% ===============================================================
\begin{thebibliography}{9}

\bibitem{gogolou2019}
A.~Gogolou, T.~Tsandilas, T.~Palpanas, and A.~Bezerianos.
\newblock Comparing Similarity Perception in Time Series Visualizations.
\newblock \emph{IEEE Transactions on Visualization and Computer Graphics},
25(1):523--533, 2019.
\newblock \url{https://doi.org/10.1109/TVCG.2018.2865077}

\bibitem{zhao2018}
J.~Zhao and L.~Itti.
\newblock shapeDTW: shape Dynamic Time Warping.
\newblock \emph{Pattern Recognition}, 74:171--184, 2018.
\newblock \url{https://arxiv.org/abs/1606.01601}

\bibitem{dtwpluss}
A.~Srivastava.
\newblock DTW+S: Shape-based Comparison of Time-series with Ordered Local Trend.
\newblock \emph{arXiv:2309.03579}, 2023 (rev.\ 2024).
\newblock \url{https://arxiv.org/abs/2309.03579}

\bibitem{trendgnn}
M.~Tian and A.~Srivastava.
\newblock TrendGNN: Towards Understanding of Epidemics, Beliefs, and Behaviors.
\newblock \emph{arXiv:2512.00421}, 2025.
\newblock \url{https://arxiv.org/abs/2512.00421}

\bibitem{esdtw2023}
ESDTW: Extrema-based shape dynamic time warping.
\newblock \emph{Expert Systems with Applications}, 2023.
\newblock \url{https://www.sciencedirect.com/science/article/abs/pii/S0957417423029342}

\bibitem{dubuc2024}
Multiscale Dubuc: A New Similarity Measure for Time Series.
\newblock \emph{arXiv:2411.10418}, 2024.
\newblock \url{https://arxiv.org/abs/2411.10418}

\bibitem{seeingeye2025}
S.~Long, A.~Chatzimparmpas, E.~Alexander, M.~Kay, and J.~Hullman.
\newblock Seeing Eye to AI? Applying Deep-Feature-Based Similarity Metrics to
Information Visualization.
\newblock \emph{ACM CHI}, 2025.
\newblock \url{https://mucollective.northwestern.edu/project/2025-perceptual-similarity-metric}

\bibitem{empirical2014}
J.~Serr\`a and J.~L.~Arcos.
\newblock An Empirical Evaluation of Similarity Measures for Time Series Classification.
\newblock \emph{arXiv:1401.3973}, 2014.
\newblock \url{https://arxiv.org/abs/1401.3973}

\bibitem{ordinal2026}
Scalable and Interpretable Representation Alignment with Ordinal Similarity
(Triplet / Quadruplet Similarity Indices).
\newblock \emph{arXiv:2606.16379}. \verify{exact authors, title, year}
\newblock \url{https://arxiv.org/abs/2606.16379}

\end{thebibliography}

\end{document}
